\section{Interpretation taxonomy: where is singular reference fixed?}
\label{sec:taxonomy}

The previous sections have now done enough work that the interpretation discussion can be stated as a
same-scope exhaustion theorem. The relevant question is not simply which interpretation one prefers.
It is this: \emph{where, if anywhere, is the work of fixing singular reference done strongly enough for
measurement records to function as theory-level evidence on the fixed scope?}

\subsection{Interpretation packages and the loci of reference-fixing work}

\begin{definition}[Interpretation package]
An \emph{interpretation package} for a measurement scenario $\mathcal{M}$ is a triple
\[
\mathfrak{I}=(T,\Lambda,\mathcal{D}^*),
\]
where:
\begin{enumerate}
    \item $T=(\Omega,\Delta)$ is the physical theory, with ontology or state-description $\Omega$ and
    dynamics $\Delta$;
    \item $\Lambda$ is the class of further resources appealed to in fixing the referents of measurement
    discourse, including added ontology, semantic restrictions, or discourse-level policies; and
    \item $\mathcal{D}^*$ is the class of diachronic measurement claims the package is intended to
    underwrite.
\end{enumerate}
\end{definition}

\begin{definition}[Discourse-committing and discourse-deflationary packages]
An interpretation package $\mathfrak{I}$ is \emph{discourse-committing} if it takes the claims in
$\mathcal{D}^*$ to be truth-apt and confirmation-guiding through the measurement records of the
scenario. It is \emph{discourse-deflationary} if it does not take the theory to underwrite those claims in
the required truth-apt diachronic sense.
\end{definition}

\begin{definition}[Dynamic selection]
A discourse-committing package realizes \emph{dynamic selection} if singular reference is fixed by the
physical evolution law itself: for the relevant measurement episode, the dynamics $\Delta$ alone
suppresses or removes outcome multiplicity strongly enough that exactly one post-measurement history
remains eligible as the referent relevant to $\mathcal{D}^*$.
\end{definition}

\begin{definition}[Ontic privileging]
A discourse-committing package realizes \emph{ontic privileging} if singular reference is fixed by
additional ontological structure not determined by the bare dynamics alone. Equivalently, the
physical description still requires an extra state variable, trajectory, configuration, or comparable
ontic fact to designate which history is the actual referent.
\end{definition}

\begin{definition}[Semantic restriction]
A discourse-committing package realizes \emph{semantic restriction} if the physical theory remains
non-selective, but admissible reference assignments are restricted strongly enough that the surviving
assignments form a single $\MReq$-class on the target discourse.
\end{definition}

\subsection{From stabilization to structural exhaustion}

\begin{lemma}[Discourse commitment forces stabilization]
\label{lem:discourse-commitment-stabilization}
If an interpretation package $\mathfrak{I}$ is discourse-committing, then it must stabilize admissible
reference assignments strongly enough to avoid the role-divergent case of
Lemma~\ref{thm:two-case-exhaustion}.
\end{lemma}

\begin{proof}
A discourse-committing package takes theory-level confirmation, memory, anticipation, deliberation,
or some combination of them to be truth-apt and guided by measurement records. By
Theorem~\ref{thm:reference-stabilization-necessity}, any such use requires admissible reference
assignments to collapse strongly enough that fixation commutes with admissible redescription on the
target discourse. So discourse commitment forces stabilization.
\end{proof}

\begin{definition}[Structural locus of singular-reference fixing]
\label{def:selector-locus}
A same-scope resource that purports to fix singular reference does its standing-bearing work at
exactly one of the following loci:
\begin{enumerate}
    \item \textbf{boundary locus}: by acting as an admissibility-relevant selector or parameter at the
    pre-admissible boundary;
    \item \textbf{quotient locus}: by refining the quotient-level standing-bearing content itself through a
    new invariant or privileged continuation relation; or
    \item \textbf{representation locus}: by restricting how already-admissible representatives or uses are
    allowed to count without changing the quotient-level physical package.
\end{enumerate}
\end{definition}

\begin{theorem}[Structural locus trichotomy]
\label{thm:selector-locus-trichotomy}
Any admissibility-preserving same-scope selector must operate at exactly one of the loci in
Definition~\ref{def:selector-locus}. No fourth same-scope locus exists.
\end{theorem}

\begin{proof}
By Proposition~\ref{prop:substrate-consequences}, any same-scope admissibility-bearing difference
must appear either in standing classification or in licensed continuation. The fixed downstream
substrate closes any second same-scope invariant except as bookkeeping or scope change, and the
role-decomposition results imported with that substrate leave no fourth admissibility-bearing role
beside anchor-level, quotient/tensor-level, and representational/skin-level work. A selector that acts
before quotient-level content is fixed therefore acts at the boundary; a selector that changes
standing-bearing content acts on the quotient; and a selector that leaves quotient content untouched
but changes how representatives may count acts via representation. If it does none of these, it is
bookkeeping only and not a selector. So no fourth same-scope locus remains.
\end{proof}

\begin{corollary}[Boundary selectors are forbidden]
\label{cor:boundary-selectors-forbidden}
If a same-scope selector operates at the boundary locus, it introduces admissibility-relevant
parameterization or selection and is therefore forbidden on the fixed scope.
\end{corollary}

\begin{proof}
By Theorem~\ref{thm:selector-locus-trichotomy}, a boundary-locus selector acts before quotient-level
content is fixed. But Theorem~\ref{thm:compressed-substrate-necessity} and
Lemma~\ref{lem:bivalence-boundary-excludes-parameterization} exclude same-scope admissibility-
relevant parameterization or selection at the boundary. So a boundary selector is forbidden.
\end{proof}

\begin{corollary}[Quotient refinements force new invariant structure]
\label{cor:quotient-refinement-scope-change}
If a same-scope selector operates at the quotient locus, then it refines standing-bearing content by a
new invariant or privileged continuation relation. On the fixed scope that is not an additional neutral
interpretation layer; it is a scope change relative to the non-selective base package.
\end{corollary}

\begin{proof}
A quotient-locus selector changes which standing-bearing content counts as the relevant continuation.
By Theorem~\ref{thm:equivalence-collapse}, no alternative same-scope admissibility-preserving
quotient survives beside $\MReq$. So any quotient refinement either coincides with the existing
quotient or introduces a new same-scope invariant. The latter is excluded by the fixed downstream
substrate except as scope change.
\end{proof}

\begin{corollary}[Representation restrictions yield semantic restriction or deflation]
\label{cor:representation-restriction-semantic-or-deflation}
If a same-scope selector operates only at the representation locus while leaving the physical package
non-selective, then either it restricts admissible representatives to one quotient class (semantic
restriction) or it refuses the target truth-apt discourse (discourse deflation).
\end{corollary}

\begin{proof}
A purely representational selector cannot change boundary status or quotient-level content without
leaving the representation locus. If the target discourse remains truth-apt and theory-relative, the only
way such a selector can matter is by restricting admissible representatives strongly enough that their
use collapses into one quotient class; that is semantic restriction. If the target discourse is instead
refused, the package is discourse-deflationary. No third representation-locus option remains.
\end{proof}

\begin{theorem}[Selector-exhaustion theorem]
\label{thm:selector-exhaustion}
Any admissibility-preserving selector must do exactly one of the following:
\begin{enumerate}
    \item introduce boundary parameterization or a boundary selector, which is forbidden;
    \item refine standing-bearing quotient content by a new invariant or privileged continuation
    relation, which is a scope change on the fixed substrate;
    \item restrict representation so that admissible assignments collapse into one quotient class,
    which is semantic restriction; or
    \item refuse the target truth-apt discourse, which is discourse deflation.
\end{enumerate}
No fifth same-scope selector route exists.
\end{theorem}

\begin{proof}
By Theorem~\ref{thm:selector-locus-trichotomy}, every same-scope selector operates at the boundary,
on the quotient, or via representation. Boundary-locus selectors are forbidden by
Corollary~\ref{cor:boundary-selectors-forbidden}. Quotient-locus selectors refine standing-bearing
content and therefore introduce a new invariant or privileged continuation relation; by
Corollary~\ref{cor:quotient-refinement-scope-change} this is a scope change on the fixed substrate.
Representation-locus selectors fall under
Corollary~\ref{cor:representation-restriction-semantic-or-deflation}: they are either semantic
restriction or discourse deflation. These cases are exhaustive, so no fifth same-scope selector route
remains.
\end{proof}

\begin{theorem}[Interpretation taxonomy]
\label{thm:selector-taxonomy}
Every interpretation package $\mathfrak{I}$ for a measurement scenario falls into exactly one of the
following structural classes:
\begin{enumerate}
    \item \textbf{Dynamic selection.} Singular reference is fixed by physical dynamics.
    \item \textbf{Ontic privileging.} Singular reference is fixed by added ontological structure or a
    privileged trajectory.
    \item \textbf{Semantic restriction.} The physical theory remains non-selective, but admissible
    reference assignments are restricted to a single standing-equivalence class.
    \item \textbf{Discourse deflation.} The package does not take the theory to underwrite the target
    class of truth-apt diachronic measurement discourse.
\end{enumerate}
No fifth structural locus remains under the fixed downstream substrate.
\end{theorem}

\begin{proof}
Take any interpretation package $\mathfrak{I}$. If it is not discourse-committing, it is discourse-
deflationary and we are in class~(4).

Assume instead that $\mathfrak{I}$ is discourse-committing. By
Lemma~\ref{lem:discourse-commitment-stabilization}, it must stabilize admissible reference
assignments strongly enough to avoid quotient-relevant divergence. If the stabilizing work is done by
changing the physical dynamics so that only one continuation remains eligible, the package realizes
dynamic selection. If the stabilizing work is done by additional ontological structure not fixed by the
bare dynamics, the package realizes ontic privileging. If the physical theory remains non-selective and
the stabilization is achieved only by restricting admissible representatives to one quotient class, the
package realizes semantic restriction. By Theorem~\ref{thm:selector-exhaustion}, no fifth same-scope
selector route exists. Hence the listed classes are exhaustive and mutually exclusive.
\end{proof}
\subsection{The comparative interpretation matrix}

\begin{table}[htbp]
\centering
\small
\begin{tabular}{>{\raggedright\arraybackslash}p{0.18\textwidth} >{\raggedright\arraybackslash}p{0.18\textwidth} >{\raggedright\arraybackslash}p{0.22\textwidth} >{\raggedright\arraybackslash}p{0.31\textwidth}}
\toprule
\textbf{Class} & \textbf{Selector locus} & \textbf{What is fixed} & \textbf{Representative strategies} \\
\midrule
Dynamic selection & Physical dynamics & One continuation by law & Collapse dynamics, objective reduction \\
Ontic privileging & Added ontology & One continuation by extra physical fact & Hidden-variable, actual-history, privileged trajectory \\
Semantic restriction & Admissible reference assignments & One standing-equivalence class of admissible assignments & Everettian semantic-fixation programs, typed successor restrictions \\
Discourse deflation & Nowhere inside the theory & The target discourse is refused & Instrumentalism, QBism, other deflationary responses \\
\bottomrule
\end{tabular}
\caption{Comparative interpretation matrix. The classification question is not which interpretation is
metaphysically best, but where the work of fixing singular reference is done if it is done at all.}
\label{tab:selector-taxonomy}
\end{table}

\subsection{Where the familiar interpretations land}

\begin{proposition}[Collapse strategies instantiate dynamic selection]
\label{prop:collapse-dynamic}
Collapse strategies such as GRW instantiate class~(1) of
Theorem~\ref{thm:selector-taxonomy}: they place the work of singular-reference fixation in the
physical dynamics itself.
\end{proposition}

\begin{proof}
A collapse theory changes the dynamics so that, for the relevant measurement episode, one outcome
history is physically singled out rather than multiple stable histories remaining equally realized.
That is exactly dynamic selection.
\end{proof}

\begin{proposition}[Hidden-variable strategies instantiate ontic privileging]
\label{prop:hidden-ontic}
Hidden-variable or actual-history strategies instantiate class~(2) of
Theorem~\ref{thm:selector-taxonomy}: they place the work of singular-reference fixation in added
ontological structure rather than in the bare dynamics alone.
\end{proposition}

\begin{proof}
Such strategies retain a multiplicity in the wavefunction-level description but add further structure ---
for example an actual configuration or privileged trajectory --- that designates which history is the
one relevant to diachronic discourse. That is ontic privileging.
\end{proof}

\begin{proposition}[Bare Everett is not yet a completed discourse-committing strategy]
\label{prop:bare-everett-not-complete}
The bare Everettian package is physically non-selective and therefore does not by itself realize class~(1),
class~(2), or class~(3). To underwrite the target discourse, it must either be completed by semantic
restriction or move to discourse deflation.
\end{proposition}

\begin{proof}
Theorem~\ref{thm:everett_nonselection} and Corollary~\ref{cor:everettian-nonselection} show that the
bare Everettian package is physically non-selective. So it does not realize dynamic selection. By
construction it adds no extra ontology that privileges one branch, so it does not realize ontic
privileging. It realizes semantic restriction only if admissible reference assignments are restricted
strongly enough to collapse into one standing-equivalence class. Bare Everett alone does not do that.
Therefore it is not yet a completed discourse-committing strategy. To underwrite the target discourse,
it must either be completed by semantic restriction or retreat to discourse deflation.
\end{proof}

\begin{corollary}[Bare Everettian exclusion of the first three loci]
\label{cor:bare-everett-exhaustion}
If one retains the bare Everettian package without adding new dynamics, without adding ontological
selection structure, and without imposing an admissible same-scope restriction on reference assignments,
then the package realizes none of the first three loci of Theorem~\ref{thm:selector-exhaustion}. Relative
to the target explanatory practice, only discourse deflation remains.
\end{corollary}

\begin{proof}
Theorem~\ref{thm:everett_nonselection} excludes dynamic selection in the bare Everettian package.
Proposition~\ref{prop:bare-everett-not-complete} excludes ontic privileging and shows that the bare
package does not by itself collapse admissible assignments into a single standing-equivalence class.
Therefore, absent additional same-scope structure, the bare package realizes none of the first three loci
of Theorem~\ref{thm:selector-exhaustion}. If the target explanatory practice is nonetheless refused, the
package moves to discourse deflation.
\end{proof}

\begin{proposition}[Everettian semantic-fixation strategies instantiate semantic restriction]
\label{prop:everett-semantic}
Any Everettian strategy that keeps the physical theory non-selective while restricting admissible
diachronic reference strongly enough for records to function as evidence falls under class~(3) of
Theorem~\ref{thm:selector-taxonomy}.
\end{proposition}

\begin{proof}
Such a strategy leaves the underlying physical package non-selective, so it is neither dynamic selection
nor ontic privileging. If it still underwrites theory-level confirmation, memory, anticipation, and
agency, then by Lemma~\ref{lem:discourse-commitment-stabilization} it must stabilize admissible
reference assignments. Since the work is not being done by a physical selector, it is being done by
restricting admissible assignments themselves. That is exactly semantic restriction.
\end{proof}

\begin{proposition}[QBist and instrumentalist strategies instantiate discourse deflation]
\label{prop:qbism-deflation}
QBist and instrumentalist strategies \citep{fuchs2014} fall under class~(4) of
Theorem~\ref{thm:selector-taxonomy} insofar as they do not take quantum theory to underwrite the
target class of truth-apt diachronic measurement discourse.
\end{proposition}

\begin{proof}
On QBist and related instrumentalist views, quantum states and outcomes are not taken to describe an
objective diachronic measurement history in the sense needed by the standing conditional. Instead they
organize expectations, betting commitments, or instrumental predictions. The package therefore does
not aim to secure singular reference for the target discourse through the theory itself. That is not a
failed stabilization strategy; it is a refusal of the target explanatory burden. Hence these views
instantiate discourse deflation.
\end{proof}

\begin{remark}[Relational and perspectival views]
Some relational or perspectival views may not advertise themselves under any of the four standard
labels \citep{rovelli1996}. The present taxonomy still applies. If such a view is discourse-committing and
physically non-selective, it must either supply semantic restriction or add physical selection. If it
instead treats the relevant diachronic discourse as non-fundamental or merely perspective-bound in a
way that refuses theory-level truth-aptness, it moves toward discourse deflation. The exact placement
of particular relational programs is therefore an application question, not a challenge to the taxonomy.
\end{remark}
